\chapter{حلقه‌ها}%
\label{ch:loops}

رایانه در یک کار بسیار توانمند است: \emph{تکرار}. کارهایی که برای انسان
خسته‌کننده و طولانی‌اند مثلاً جمع‌زدنِ هزار عدد یا چاپِ صد سطر را رایانه
در یک چشم‌به‌هم‌زدن انجام می‌دهد. ابزارِ این تکرار، \textbf{حلقه} است. در این
فصل با همهٔ حلقه‌های سلام آشنا می‌شویم.

\section{تاوقتی: تا زمانی که شرط برقرار است}

ساده‌ترین حلقه، \code{تاوقتی} (در انگلیسی \code{until}) است. توجه کنید که
سلام هیچ کلیدواژه‌ای به نامِ \code{while} ندارد؛ \code{تاوقتی} تنها حلقهٔ
شرط‌محورِ زبان است. این حلقه بدنهٔ خود را \emph{تا زمانی که} یک شرط درست
است، بارها و بارها اجرا می‌کند. به این مثالِ واقعی نگاه کنید که اعدادِ ۱ تا
۱۰ را جمع می‌زند:

\begin{salam}
تابع اصلی:
    متغیر مجموع : صحیح = 0
    متغیر i : صحیح = 1
    تاوقتی i <= 10:
        مجموع = مجموع + i
        i = i + 1
    تمام
    چاپ("مجموع ۱ تا ۱۰ برابر است با", مجموع)
تمام
\end{salam}

این حلقه را گام‌به‌گام دنبال کنیم:
\begin{enumerate}
  \item نخست \code{مجموع} برابرِ ۰ و \code{i} برابرِ ۱ است.
  \item سلام شرطِ \code{i <= 10} را می‌آزماید. درست است، پس بدنه اجرا می‌شود:
        \code{مجموع} برابرِ ۱ و \code{i} برابرِ ۲ می‌شود.
  \item دوباره شرط آزموده می‌شود. باز درست است... و این کار تا وقتی \code{i}
        به ۱۱ برسد ادامه می‌یابد.
  \item وقتی \code{i} برابرِ ۱۱ شد، شرط \emph{نادرست} می‌شود و حلقه پایان
        می‌یابد.
\end{enumerate}

خروجی: \code{مجموع ۱ تا ۱۰ برابر است با 55}.

\begin{warn}
هر حلقهٔ \code{تاوقتی} به دو چیز نیاز دارد: شرطی که سرانجام نادرست شود، و
تغییری در بدنه که ما را به آن نزدیک کند. اگر خطِ \code{i = i + 1} را فراموش
کنید، \code{i} همیشه ۱ می‌ماند، شرط همیشه درست است و حلقه \emph{هرگز پایان
نمی‌یابد}! به این اشتباهِ رایج، \textbf{حلقهٔ بی‌پایان} می‌گویند.
\end{warn}

\section{تکرار: حلقهٔ شمارشی}

بسیار پیش می‌آید که می‌خواهیم کاری را به تعدادِ مشخصی تکرار کنیم، گاهی هم با
یک شمارنده. حلقهٔ \code{تکرار} (معادلِ \code{repeat}) دقیقاً برای همین ساخته
شده است. ساده‌ترین شکلِ آن فقط عددِ تکرار را می‌گیرد و شمارنده‌ای در اختیار
نمی‌گذارد:

\begin{salam}
تابع اصلی:
    تکرار 5:
        چاپ "سلام"
    تمام
تمام
\end{salam}

این برنامه پنج‌بار «سلام» چاپ می‌کند. اما وقتی به شمارهٔ هر تکرار هم نیاز
داریم مثلاً برای چاپِ جدولِ ضرب پس از عددِ تکرار کلمهٔ \code{با} و یک نام
می‌آوریم؛ سلام آن نام را از \code{۰} تا یکی کمتر از عددِ تکرار مقداردهی
می‌کند:

\begin{salam}
تابع اصلی:
    تکرار 5 با i:
        چاپ("شمارهٔ", i)
    تمام
تمام
\end{salam}

خروجیِ این حلقه شمارهٔ \code{۰} تا \code{۴} است (نه ۱ تا ۵)؛ شمارش همیشه از
صفر آغاز می‌شود.

\begin{warn}
هرگز برای شمردنِ \code{۰} تا \code{n-1} ننویسید \code{تکرار 0 تا n-1 با i:}.
وقتی \code{n} برابرِ صفر باشد، این نوشتار به \code{0 تا -1} تبدیل می‌شود که
\emph{رو به پایین} می‌شمارد؛ یک اشتباهِ رایج و موذی! برای شمارشِ \code{۰} تا
\code{n-1} همیشه همان شکلِ ساده را به‌کار ببرید: \code{تکرار n با i:}.
\end{warn}

اگر بخواهیم شمارنده از عددی دلخواه (نه لزوماً صفر) آغاز شود و تا عددی دیگر
پیش برود، شکلِ \emph{بازه‌ایِ} حلقهٔ \code{تکرار} به کار می‌آید:
\code{تکرار آ تا ب با i:} مقدارِ \code{i} را از \code{آ} تا \code{ب} (شاملِ
خودِ \code{ب}) پیش می‌برد. این مثالِ واقعی جدولِ ضربِ ۵ را چاپ می‌کند:

\begin{salam}
تابع اصلی:
    تکرار 1 تا 10 با i:
        چاپ(5, "*", i, "=", 5 * i)
    تمام
تمام
\end{salam}

خطِ \code{تکرار 1 تا 10 با i:} چنین می‌خواند: «از ۱ تا ۱۰ (هر دو شامل)، این
بدنه را اجرا کن و در هر بار مقدارِ جاری را در \code{i} بگذار».

اگر بخواهیم شمارنده با گامی غیر از ۱ پیش برود، پس از پایانِ بازه کلمهٔ
\code{گام} و عددِ گام را می‌آوریم:

\begin{salam}
تابع اصلی:
    تکرار 0 تا 100 گام 10 با i:
        چاپ i
    تمام
تمام
\end{salam}

این حلقه اعدادِ \code{۰، ۱۰، ۲۰، ...، ۱۰۰} را چاپ می‌کند. جهتِ شمارش خودکار
است: اگر آغازِ بازه از پایانش کوچک‌تر باشد، شمارنده بالا می‌رود؛ در
غیرِاین‌صورت پایین می‌آید.

هر پنج شکلِ \code{تکرار} را می‌توان این‌گونه خلاصه کرد:

\begin{itemize}
  \item \code{تکرار N:} بدنه را \code{N} بار اجرا می‌کند، بدونِ شمارنده.
  \item \code{تکرار N با i:} همان کار را می‌کند، و \code{i} را از \code{۰}
        تا \code{N-1} مقداردهی می‌کند.
  \item \code{تکرار آ تا ب:} برای هر مقدار از \code{آ} تا \code{ب} (شامل) یک
        بار اجرا می‌شود، بدونِ شمارنده.
  \item \code{تکرار آ تا ب با i:} همان بازه را می‌پیماید و مقدارِ جاری را در
        \code{i} می‌گذارد.
  \item \code{تکرار آ تا ب گام گ با i:} مانندِ بالا، اما با گام‌هایی به
        اندازهٔ \code{گ}.
\end{itemize}

\begin{tip}
نامِ \code{i} برای شمارندهٔ حلقه یک سنّتِ دیرینه است (کوتاه‌شدهٔ \emph{index}).
برای حلقه‌های تودرتو از \code{i}، \code{j} و \code{k} استفاده می‌کنند. البته
می‌توانید نام‌های فارسیِ گویاتر مانندِ \code{ردیف} و \code{ستون} هم به‌کار
ببرید.
\end{tip}

\begin{note}
\code{تکرار} برای شمردن ساخته شده، نه برای گشتن در عناصرِ یک آرایه. وقتی
مسئله واقعاً «برای هر عنصرِ این مجموعه...» است نه «N بار» یا «از آ تا ب»
حلقهٔ \code{هر} (معادلِ \code{each}) انتخابِ درست‌تری است؛ چون مستقیماً روی
عناصر می‌چرخد و نیازی به شمارش دستیِ اندیس‌ها ندارید:

\begin{salam}
تابع اصلی:
    نام‌ها := ["آرش", "بیتا", "سینا"]
    هر ن در نام‌ها:
        چاپ ن
    تمام
تمام
\end{salam}

دربارهٔ آرایه‌ها و حلقهٔ \code{هر} در فصلِ آرایه‌ها بیشتر خواهیم خواند.
\end{note}

\section{شکل‌کشی با حلقه‌های تودرتو}

وقتی یک حلقه را \emph{درونِ} حلقهٔ دیگر بگذاریم، کارهای جالبی می‌توان کرد. این
مثالِ واقعی یک مثلثِ ستاره می‌کشد نمونه‌ای محبوب نزدِ تازه‌کارها:

\begin{salam}
// مثلث ستاره
تابع اصلی:
    تکرار 1 تا 5 با ردیف:
        متغیر خط : رشته = ""
        تکرار ردیف با ستون:
            خط = خط + "*"
        تمام
        چاپ خط
    تمام
تمام
\end{salam}

خروجی:

\begin{salamen}[language={}]
*
**
***
****
*****
\end{salamen}

حلقهٔ بیرونی روی ردیف‌ها (۱ تا ۵) می‌چرخد، و برای هر ردیف، حلقهٔ درونی به تعدادِ
شمارهٔ آن ردیف ستاره به رشتهٔ \code{خط} می‌افزاید. سپس خط چاپ می‌شود. این الگوی
«ساختنِ یک رشته در حلقهٔ درونی و چاپِ آن در حلقهٔ بیرونی» را بارها به‌کار خواهید
برد.

\begin{tip}
همین برنامه را می‌توان با جدولِ ضرب هم نوشت: کافی است حلقهٔ درونی را
\code{تکرار 1 تا ردیف با ستون:} بنویسید و به‌جای افزودنِ ستاره، حاصلِ
\code{ردیف * ستون} را به خط بیفزایید. مثال‌های واقعیِ سلام چنین برنامه‌ای
نیز دارند.
\end{tip}

\section{بشکن و ادامه: کنترلِ دقیق‌ترِ حلقه}

گاهی می‌خواهیم پیش از پایانِ طبیعیِ حلقه از آن بیرون بیاییم، یا یک تکرار را رد
کنیم. برای این دو کار، \code{بشکن} (معادلِ \code{break}) و \code{ادامه}
(معادلِ \code{continue}) را داریم. به این مثالِ واقعی نگاه کنید:

\begin{salam}
تابع اصلی:
    متغیر مجموع : صحیح = 0
    متغیر i : صحیح = 0
    تاوقتی i < 20:
        i = i + 1
        اگر i % 2 == 0:
            ادامه
        تمام
        اگر i > 9:
            بشکن
        تمام
        مجموع = مجموع + i
    تمام
    چاپ مجموع
تمام
\end{salam}

در این حلقه:
\begin{itemize}
  \item \code{ادامه} وقتی \code{i} زوج است اجرا می‌شود و باعث می‌شود سلام
        \emph{باقیِ} بدنه را رد کند و به تکرارِ بعد برود. پس اعدادِ زوج هرگز به
        مجموع افزوده نمی‌شوند.
  \item \code{بشکن} وقتی \code{i} از ۹ بزرگ‌تر شود، حلقه را \emph{به‌کلی}
        پایان می‌دهد.
\end{itemize}
در نتیجه فقط اعدادِ فردِ ۱ تا ۹ (یعنی ۱، ۳، ۵، ۷، ۹) جمع می‌شوند و خروجی
\code{25} می‌شود.

\begin{tip}
\code{بشکن} برای زمانی عالی است که دنبالِ چیزی می‌گردید و همین که یافتیدش، دیگر
نیازی به ادامهٔ جست‌وجو نیست. \code{ادامه} هم برای رد کردنِ موارد ناخواسته
(مثلِ اعدادِ زوج در مثالِ بالا) به‌کار می‌آید. هر دو در حلقه‌های \code{تاوقتی}،
\code{تکرار} و \code{هر} یکسان کار می‌کنند.
\end{tip}

\section{مثالِ کامل: بازیِ فیزباز}

«فیزباز» تمرینی کلاسیک است که حلقه، شرط و باقی‌مانده را با هم به کار می‌گیرد.
قاعده ساده است: از ۱ تا ۱۵ بشمار؛ به‌جای مضرب‌های ۳ بنویس «فیز»، به‌جای
مضرب‌های ۵ بنویس «باز»، و به‌جای مضرب‌های هر دو بنویس «فیزباز». این مثالِ واقعیِ
سلام است:

\begin{salam}
// فیزباز
تابع اصلی:
    تکرار 1 تا 15 با i:
        اگر i % 15 == 0:      چاپ "فیزباز"
        وگرنه i % 3 == 0: چاپ "فیز"
        وگرنه i % 5 == 0: چاپ "باز"
        وگرنه:               چاپ i
        تمام
    تمام
تمام
\end{salam}

\begin{warn}
چرا نخست \code{i \% 15 == 0} را می‌آزماییم؟ چون عددی که هم بر ۳ و هم بر ۵
بخش‌پذیر است، حتماً بر ۱۵ هم بخش‌پذیر است. اگر این شرط را \emph{آخر}
می‌گذاشتیم، عددِ ۱۵ زودتر در شرطِ \code{\% 3} گیر می‌افتاد و به‌اشتباه فقط
«فیز» چاپ می‌شد. باز هم می‌بینیم که \emph{ترتیبِ شرط‌ها} سرنوشت‌ساز است.
\end{warn}

\section{خلاصهٔ فصل}

\begin{itemize}
  \item \code{تاوقتی} تا زمانی که شرط درست است تکرار می‌کند؛ تنها حلقهٔ
        شرط‌محورِ سلام است.
  \item \code{تکرار} حلقهٔ شمارشی است: \code{تکرار N:}، \code{تکرار N با i:}،
        \code{تکرار آ تا ب:}، \code{تکرار آ تا ب با i:} و \code{تکرار آ تا ب
        گام گ با i:}.
  \item \code{هر عنصر در مجموعه:} برای گشتن در عناصرِ یک آرایه یا مجموعه
        به‌کار می‌رود، نه برای شمردن.
  \item حلقه‌های تودرتو برای شکل‌کشی و جدول‌ها عالی‌اند.
  \item \code{بشکن} حلقه را پایان می‌دهد و \code{ادامه} به تکرارِ بعد می‌پرد.
  \item مراقبِ حلقه‌های بی‌پایان، و مراقبِ \code{تکرار 0 تا n-1} باشید!
\end{itemize}

\section{تمرین‌ها}

\exercise با حلقهٔ \code{تکرار} (شکلِ بازه‌ای)، اعدادِ ۱ تا ۲۰ را چاپ کنید.

\exercise با حلقهٔ \code{تاوقتی}، جمعِ اعدادِ زوجِ ۱ تا ۱۰۰ را حساب کنید.

\exercise با حلقهٔ \code{تکرار} (شکلِ بازه‌ای با \code{با i})، فاکتوریلِ عددِ
۵ را حساب کنید (یعنی $1\times2\times3\times4\times5$).

\exercise یک مثلثِ \emph{راست‌چین} با ۵ ردیف بکشید (هر ردیف چند فاصله و سپس
ستاره دارد).

\exercise جدولِ ضربِ کاملِ ۱ تا ۵ را با دو حلقهٔ تودرتو چاپ کنید.

\exercise با \code{بشکن}، برنامه‌ای بنویسید که اعداد را از ۱ بشمارد و به محضِ
رسیدن به نخستین عددی که بر ۷ بخش‌پذیر است، آن را چاپ کند و بایستد.
